\documentclass{article}
\usepackage[main]{neurips_2026}
\usepackage[utf8]{inputenc}
\usepackage[T1]{fontenc}
\usepackage{courier}
\usepackage{hyperref}
\usepackage{url}
\usepackage{booktabs}
\usepackage{amsfonts}
\usepackage{amsmath}
\usepackage{amssymb}
\usepackage{graphicx}
\usepackage{longtable}
\usepackage{xcolor}
\input{manuscript/claim_values}
\title{Reliability Beyond Marginal Coverage for Observed Compositions}
\author{Anonymous}
\begin{document}
\maketitle
\begin{abstract}
Conformal regions for an observed composition require simultaneous-vector coverage, subgroup diagnostics and a common measure of efficiency. We compare split-conformal scores and calibration schemes in controlled logistic-normal regression using 100 independent full-dataset repetitions per main cell, three or ten components, and a shared fitted predictor. At three components with heterogeneous noise and 499 calibration observations, global Aitchison joint coverage is $\GlobalJoint{}$, but high-noise-group coverage is only $\GlobalHigh{}$. A shift increasing that group's frequency reduces marginal coverage to $\GlobalShift{}$, while observed-group calibration reaches $\GroupShift{}$. Learned scaling improves adaptation but still misses the group target in a ten-component case; correct oracle weighting restores shifted marginal coverage without guaranteeing group coverage. Common Euclidean simplex-area estimates reveal setting-dependent efficiency differences, with coverage reported alongside size. Coordinatewise calibration can substantially undercover the whole vector, and zero replacement requires distinguishing direct raw-response membership from an explicit inverse-image region, which inherits transformed coverage for original responses. The contribution is a reproducible diagnostic evaluation under explicit synthetic mechanisms, not a new coverage theorem or a universal method ranking.
\end{abstract}
\input{manuscript/sections/introduction}
\input{manuscript/sections/related_work}
\input{manuscript/sections/methods}
\input{manuscript/sections/results}
\input{manuscript/sections/limitations}
\clearpage
\bibliographystyle{plainnat}
\bibliography{manuscript/references}
\clearpage
\appendix
\input{manuscript/sections/theory}
\input{manuscript/sections/details}
\input{manuscript/sections/workflow}
\input{manuscript/sections/reproduction}
\clearpage
\input{manuscript/sections/displays}
\clearpage
\input{manuscript/checklist}
\end{document}
